\newpage

\section{Разведочный анализ данных}

\subsection{Общая статистика и пропуски}

Датасет A содержит 5739 записей. Анализ пропущенных значений (Рис.~\ref{fig:missing}) показал, что наибольшая доля пропусков наблюдается в полях \texttt{cb\_url} (56.9\%), \texttt{linkedin\_url} (54.7\%) и \texttt{year\_founded} (44.9\%). Поля \texttt{country} и \texttt{location} имеют около 29\% пропусков, \texttt{num\_founders} --- 26.8\%. Основные поля из датасета 2025 года (\texttt{team\_size}, \texttt{industry}, \texttt{status}) заполнены практически полностью.


\begin{figure}[htbp]
    \centering
    \includegraphics[width=1\textwidth]{missing_values.png}
    \caption{Доля пропущенных значений в Датасете A}
    \label{fig:missing}
\end{figure}

Высокая доля пропусков в полях \texttt{cb\_url}, \texttt{linkedin\_url}, \texttt{year\_founded}, \texttt{country} объясняется тем, что эти колонки были добавлены из датасета 2023 года и доступны только для 4202 компаний из 5739.

\subsection{Распределение целевой переменной}

Распределение статусов в Датасете A представлено на Рис.~\ref{fig:status}. Подавляющее большинство компаний имеют статус active --- 3961 (69.0\%). Статус inactive имеют 1011 компаний (17.6\%), acquired --- 744 (13.0\%), public --- 23 (0.4\%).

\begin{figure}[htbp]
    \centering
    \includegraphics[width=0.9\textwidth]{status_distribution.png}
    \caption{Распределение целевой переменной status в Датасете A}
    \label{fig:status}
\end{figure}

После бинаризации для сравнительного анализа success vs failure было отобрано 1778 компаний: 767 успешных (acquired или public) и 1011 неуспешных (inactive). Остальные 3961 компания со статусом active были исключены из сравнительного анализа.

\subsection{Числовые признаки}

Распределения числовых признаков представлены на Рис.~\ref{fig:numerical}. Размер команды (\texttt{team\_size}) имеет сильно правоскошенное распределение: большинство стартапов имеют небольшие команды, но встречаются компании с сотнями сотрудников. Год основания (\texttt{year\_founded}) концентрируется в диапазоне 2015--2024 годов, что отражает рост числа стартапов-выпускников акселератора Y Combinator в последнее десятилетие.

\begin{landscape}
\begin{figure}[htbp]
    \centering
    \includegraphics[width=1.5\textwidth]{numerical_distributions.png}
    \caption{Распределения числовых признаков в Датасете A}
    \label{fig:numerical}
\end{figure}
\end{landscape}

\subsection{Категориальные признаки}

\subsubsection{Батчи по годам}

Динамика количества стартапов по году батча (Рис.~\ref{fig:batch}) показывает устойчивый рост числа принимаемых в Y Combinator компаний вплоть до 2021 года, и незначительное снижение и стабилизацию начиная с 2022. При этом если в 2005--2010 годах в каждом батче выпускалось несколько десятков стартапов, то к 2021--2024 годам это число выросло до нескольких сотен. Пик же количества выпускников приходится на 2021--2022 годы.

\begin{figure}[htbp]
    \centering
    \includegraphics[width=0.9\textwidth]{batch_by_year.png}
    \caption{Количество стартапов YC по году батча}
    \label{fig:batch}
\end{figure}

\subsubsection{Топ-10 отраслей}

Наиболее представленной отраслью является B2B (Рис.~\ref{fig:industries}), за которой следуют Healthcare, Fintech и Consumer. Это отражает общую тенденцию венчурного рынка, где B2B-решения традиционно привлекают наибольшее внимание акселераторов.

\begin{figure}[htbp]
    \centering
    \includegraphics[width=0.7\textwidth]{top_industries.png}
    \caption{Топ-10 отраслей стартапов YC}
    \label{fig:industries}
\end{figure}

\subsubsection{Топ-10 стран}

Географическое распределение стартапов (Рис.~\ref{fig:countries}) показывает доминирование США, что ожидаемо для американского акселератора. Среди других стран выделяются Индия, Канада, Великобритания. Также выделяется Сингапур, который, являясь городом-государством, вошёл в топ-10 стран по количеству стартапов-выпускников YC.

\newpage

\begin{figure}[htbp]
    \centering
    \includegraphics[width=0.7\textwidth]{top_countries.png}
    \caption{Топ-10 стран по количеству стартапов YC}
    \label{fig:countries}
\end{figure}

\subsubsection{Топ-20 тегов}

Анализ тегов (Рис.~\ref{fig:tags}) показывает, что наиболее популярными направлениями являются всё связанное с искусственным интеллектом и машинным обучением, B2B, SaaS, а также отрасли, связанные с финансами и здравоохранением.

\begin{figure}[htbp]
    \centering
    \includegraphics[width=0.9\textwidth]{top_tags.png}
    \caption{Топ-20 тегов стартапов YC}
    \label{fig:tags}
\end{figure}
